\subsection*{给逻辑购买选择一个生效位置}

把分散在多次尝试里的步骤当成一项购买，需要一个判断依据：这项购买何时已经成为数据库事实？P1 在事务内生成时，随后还可以整体回滚；流水插入完成时，也仍不能供其他普通连接恢复。提交使身份、资金变化、订单、流水和原结果共同成立。在本章限定的执行规则下，可以把这一数据库提交选为逻辑购买的线性化点。

线性化要求一段并发历史能够解释为某个逐项执行的历史，每项已完成操作仿佛在调用与取得结果之间的一个点生效，同时尊重已经完成再开始的实际先后。它的用途是让调用方用单项购买的规则理解并发结果。以 S1 为例，即使内部经历多条 SQL，外部只需判断在选定点以前尚未形成该购买，在点以后已经形成完整终态。提交固定业务效果在历史中的位置；客户端随后通过成功响应、查询或结果重放获知这一事实。

这里必须先说明调用与结果属于哪个对象。逻辑 S1 从第一次发送以前保存身份并开始提交，到调用方最终取得其终态为止，可能横跨 R1、R2、R3。R1 超时或收到 503 暂时不可用，只结束这次通信尝试，逻辑购买仍可处于待确认。把每个 HTTP 尝试都解释成一次全新的购买，会把后来的结果重放误数为第二次成交；把 R2 的发送时刻选成原购买生效点，又会把已经提交、可被别人观察的 P1 错移到未来。

Raft 的客户端恢复也面临提交与获知分离的问题。Ongaro 与 Ousterhout 在讨论客户端交互时指出，即使日志已经提交，响应前的崩溃仍可使客户端重发同一命令；服务器还需要识别重复并保存原响应\cite{raft,ongaro}。这里引用这项分析，是为了区分底层提交与调用恢复的职责。本章的提交顺序与持久性由单一数据库提供，请求身份与原结果由购买协议补全。\footnote{Raft 原论文为 Ongaro 与 Ousterhout 的 \textit{In Search of an Understandable Consensus Algorithm}（2014）；Ongaro 的博士论文 \textit{Consensus: Bridging Theory and Practice}（2014）进一步分析客户端会话、重复命令与已保存响应。日志负责在副本间保留命令次序，状态机按顺序应用命令；重复过滤仍须由命令身份与结果记录完成。本章采用其中的客户端交互分析，实际验证对象为单一 PostgreSQL 数据库上的购买协议。}

\subsection*{已经生效仍可能一直未被获知}

重新从 10000、空记录开始。R1 进入事务，数据库提交 P1、O1、D1，余额为 7000；进程在返回以前退出。此后另一项明确确认的 S2 扣 3000，提交 P2，余额为 4000。R2 恢复 S1，读回 P1 的当次 7000。按提交次序解释，S1 在 S2 以前生效，两笔各自保留自己的结果，整个顺序符合购买规则。

如果因为 Alice 迟至 R2 才知道 S1 成功，就把 S1 排在 S2 以后，P1 的当次 7000 便与该次序矛盾：S2 先扣后应为 7000，S1 再扣应为 4000。获知时间晚不改变已经保存的资金顺序。成功响应经过重放才到达，仍是对原位置的报告。

\begin{figure}[htbp]
\centering
\begin{tikzpicture}
\draw[bookarrow] (0,0) -- (11.8,0);
\node[booknote,anchor=east] at (0,1.8) {数据库};
\node[bookresult,text width=2.8cm] (one) at (2,1.8) {S1 提交\\D1 当次余额 7000};
\node[bookresult,text width=2.8cm] (two) at (6,1.8) {S2 提交\\D2 当次余额 4000};
\draw[bookarrow] (one) -- (two);
\draw[dashed,draw=Muted] (2,1.1) -- (2,0);
\draw[dashed,draw=Muted] (6,1.1) -- (6,0);
\node[booknode,text width=3cm] (learn) at (9.8,-1.2) {R2 恢复 S1\\客户端才获知原 7000};
\draw[bookarrow] (one.south) -- (2,-1.2) -- (learn.west);
\node[booknote] at (5,-.6) {R1 响应丢失；原结果继续保留};
\end{tikzpicture}
\caption{提交次序固定资金历史，较晚获知仍报告原位置的结果}
\label{fig:ch06-commit-learning}
\end{figure}


这也解释了没有完整响应的操作怎样进入历史。S1 的调用尚待确认时，数据库可能已经把它提交。对整段业务历史进行解释，应允许这项尚未返回的逻辑操作已经生效；只要随后成功查询或重放，都能在同一位置接回。若它在提交前回滚，历史里就没有这次尝试的购买效果，后来恢复第一次提交时才选定生效点。

如果中断发生在 COMMIT 发出以后、应用取得确认以前，恢复者需要保留三个分支。数据库已经提交，原 \(S\) 插入零行，下一条 \code{SELECT} 返回 P1；数据库回滚，恢复可能实际插入身份，才首次完成购买；原事务还没结束，唯一检查继续等待或本次暂时返回 503。三个分支使用同一套身份步骤，服务按数据库结局继续。客户端看到 \code{unknown} 时，数据库中仍可能已有成功终态。

\begin{figure}[htbp]
\centering
\begin{tikzpicture}
\node[booknode,text width=9.5cm] (lost) at (5.5,3.5) {COMMIT 已发出，应用尚未得到确认；恢复继续使用原 S};
\node[bookresult,text width=3cm] (yes) at (1.6,1.4) {数据库已提交\\插入零行\\新 SELECT 读 P1};
\node[booknode,text width=3cm] (no) at (5.5,1.4) {数据库已回滚\\可能取得身份\\首次完成购买};
\node[booknode,text width=3cm] (wait) at (9.4,1.4) {原事务未结束\\唯一检查等待\\或本次返回 503};
\draw[bookarrow] (lost) -- (yes);
\draw[bookarrow] (lost) -- (no);
\draw[bookarrow] (lost) -- (wait);
\node[booknote] at (5.5,-.1) {连接异常对应观察不完整；三个数据库分支由相同身份步骤处理};
\end{tikzpicture}
\caption{提交确认未到达时保留三种数据库结局}
\label{fig:ch06-commit-unknown}
\end{figure}


若服务已经读回 P1，发送响应时又失败，下一次恢复仍读取 P1。这一重复读取不再有新的资金线性化点；每次读取本身有自己的观察位置，但所报告的购买事实来自同一次提交。原结果保持稳定，既给出终态，也给出这段历史中的编号与当次余额证据。

\subsection*{查询读到的是哪一个时点}

查询 \(S\) 的只读操作也需要观察位置。入口经过认证后，在权威数据库发出一条 \code{SELECT}，使用语句开始时的快照。若提交已经在这条查询开始以前完成，且记录仍保留、查询条件没有额外过滤，就应读到原终态。查询期间新发生的提交，可以暂时不在该快照中；查不到描述的是这个读取时点的事实。

设 Q1 开始时 R1 尚未提交，Q1 得到空，随后 R1 提交，再随后 Q1 的 404 到达客户端。Q1 可以放在 R1 提交以前解释，因为它与 R1 的提交过程重叠。404 到达较晚，没有使快照自动变晚。因此即使客户端刚看到 404，R1 也可能已经成功。若 Q2 在 R1 明确成功返回以后才开始，却从落后副本读到空，就无法继续声称查询反映当前权威事实；读取路径已经破坏前提。

\begin{figure}[htbp]
\centering
\begin{tikzpicture}
\draw[bookarrow] (0,0) -- (11.6,0);
\node[booknode,text width=2.8cm] (query) at (1.7,1.3) {Q1 开始 SELECT\\快照中无 P1};
\node[bookresult,text width=2.8cm] (commit) at (5.5,1.3) {R1 完整提交\\P1 已成立};
\node[booknode,text width=2.8cm] (reply) at (9.3,1.3) {Q1 的 404 到达\\报告原读取时点};
\draw[bookarrow] (query) -- (commit);
\draw[bookarrow] (commit) -- (reply);
\foreach \x in {1.7,5.5,9.3}{\draw[dashed,draw=Muted] (\x,.7) -- (\x,0);}
\node[booknote,text width=10.5cm] at (5.5,-.9) {404 到达时 P1 可以已经提交；后续查询仍按完整 S 读取\\若明确成功返回后才开始查询，应在保留的权威记录中读到终态};
\end{tikzpicture}
\caption{查询与购买重叠时，查空快照可以早于购买提交}
\label{fig:ch06-query-snapshot}
\end{figure}


同 \(S\) 异参的 POST 也没有第二次购买生效点。它找到旧绑定，在比较以后返回 409，说明本次输入与该身份已经成立的含义冲突。这个响应结束当前不符合原绑定的修改尝试，原逻辑购买的终态保持不变。若原事务还在执行，它先在身份处等待；原事务提交才取得比较依据，原事务回滚则可能由当前输入首次取得身份。服务检查的是已提交绑定，调用方仍须确保同一操作的所有尝试来自同一份确认输入。

原记录的查询范围必须与仲裁范围相同：完整四列、相同权威数据库、终态仍保留、不按订单现状排除记录。若读取遗漏会话 E1，可能拿 E2 的 K1 回答；若只查未关闭订单，原 O1 关闭以后会使 S1 恢复查空。这样的空查询来自定位错误，不构成原购买从未发生的证据。

\subsection*{余额不足为什么需要单独核对}

成功扣款的条件 \code{UPDATE} 修改账户行，其他修改者要等本事务结束。事务读到扣后 7000，随后一直持有相关修改保护直到提交；另一购买要等它提交或回滚，才能基于后续余额处理。因此在成功提交点，3000 的支出与当次 7000 可以放进同一逐笔历史。

余额不足的 \code{UPDATE} 返回零行，通常没有把一条账户修改锁留给拒绝分支。如果只有本章的购买扣款，余额只会保持或减少。某时读到余额小于 3000，到拒绝提交时它仍小于 3000；把拒绝放在提交点，也符合当时条件。拒绝没有扣款，后续账户变化不会把该终态重新开启。

\begin{figure}[htbp]
\centering
\begin{tikzpicture}
\node[booknode,text width=3.2cm] (judge) at (1.8,1.2) {条件 UPDATE 判定\\余额 \(b<3000\)};
\node[booknode,text width=3.2cm] (middle) at (5.8,1.2) {只有其他购买扣款\\余额保持或减少};
\node[bookresult,text width=3.2cm] (commit) at (9.8,1.2) {拒绝终态提交\\此时仍小于 3000};
\draw[bookarrow] (judge) -- (middle);
\draw[bookarrow] (middle) -- (commit);
\node[booknote,text width=10.5cm] at (5.8,-.2) {余额单调不增，使判断时成立的不足条件延续到提交点\\该理由使用了允许资金动作的范围};
\end{tikzpicture}
\caption{只有购买支出时，余额不足条件延续到拒绝提交}
\label{fig:ch06-rejection-monotonic}
\end{figure}


一旦加入充值，这段理由就要改变。初始 2000，A 的条件 \code{UPDATE} 得到零行；B 充值 4000 并提交到 6000；A 最后才提交余额不足。拒绝在判定时有依据，在提交时余额却已经足够。如果仍要求所有终态统一在提交处满足顺序规则，可以让首次业务判断先取得账户行的排他保护，保持到终态提交，充值也遵守相同保护。或者为拒绝选择实际判定时点，重新论证身份发布、查询与其他动作的历史。两种安排都需要额外规则，本章验证不包含并发充值。

\begin{figure}[htbp]
\centering
\begin{tikzpicture}
\node[booknode,text width=3.1cm] (a) at (1.7,2) {1. A 判断不足\\初始余额 2000};
\node[bookresult,text width=3.1cm] (b) at (5.7,2) {2. B 充值 4000\\提交到 6000};
\node[booknode,text width=3.1cm] (c) at (9.7,2) {3. A 提交拒绝\\提交时余额已足够};
\draw[bookarrow] (a) -- (b);
\draw[bookalert] (b) -- (c);
\node[booknode,text width=10.5cm] at (5.7,-.1) {扩展若仍统一采用提交点：首次判断先保护账户行，\\所有充值与购买遵守同一保护，保持到各自终态提交};
\end{tikzpicture}
\caption{并发充值打断不足条件的延续，需要另行设计与证明}
\label{fig:ch06-credit-counterexample}
\end{figure}


把这项前提明写出来，也限定了 PostgreSQL Read Committed 的作用。它提供每条语句的可见范围和单行更新冲突处理；本章的线性化说明还依赖具体受理步骤、账户更新和允许动作集合。换成跨行累计限额或外部支付，必须重新选择相应的观察与提交关系。
\FloatBarrier
